\documentclass[12pt,hyperref=unicode]{beamer}

\usepackage[utf8]{inputenc}
\usepackage[turkish]{babel}
\usepackage[T1]{fontenc}

\usepackage{amsmath}
\usepackage{amssymb}
\usepackage{amsthm}
\usepackage{enumerate}
\usepackage{tikz}
\usepackage{transparent}
\usepackage{xcolor}
\usepackage{listings}
\usepackage{verbatim}
\usepackage{graphicx}

\usepackage{pgfplots}

\usetheme{Warsaw}
\usecolortheme{default}
\definecolor{nigdeyesili_acik}{RGB}{3, 150, 166}
\definecolor{nigdeyesili_koyu}{RGB}{143, 209, 217}
\setbeamercolor{structure}{fg=nigdeyesili_acik}

\author[Dr. Mehmet CANEVİ]{Arş.~Gör.~Dr.~M.~Canevi\inst{1}}

\institute{
    \inst{1}%
    Bilgisayar Mühendisliği\\
    Mühendislik Fakültesi
}
    
\date[2025] {Ders Notları, Ocak 2025}
\setbeamertemplate{footline}[frame number]

\logo{\transparent{0.4}\includegraphics[width = 20mm]{logo}}

\definecolor{mGreen}{rgb}{0,0.6,0}
\definecolor{mGray}{rgb}{0.5,0.5,0.5}
\definecolor{mPurple}{rgb}{0.58,0,0.82}
\definecolor{backgroundColour}{rgb}{0.95,0.95,0.92}

\lstloadlanguages{C}
\lstdefinestyle{CStyle}{
    backgroundcolor=\color{backgroundColour},   
    commentstyle=\color{mGreen},
    keywordstyle=\color{red},
    numberstyle=\tiny\color{mGray},
    stringstyle=\color{mGreen},
    basicstyle=\footnotesize,
    breakatwhitespace=false,         
    breaklines=true,                 
    captionpos=b,                    
    keepspaces=true,                 
    numbers=left,                    
    numbersep=2pt,                  
    showspaces=false,                
    showstringspaces=false,
    showtabs=false,                  
    tabsize=1,
    language=C
}
\lstset{style=CStyle}

\AtBeginDocument{\shorthandoff{=}}
\title[Ders 5] {Diziler II}
\begin{document}
%%%%%%%%%%%%%%%%%%%%%%%%%%%%%%%%%%%%%%%%%%%%%%%%%%%%%%%%%%%%%%%%%%%%%%%%%%%%%%%%
\frame{\titlepage}
\begin{frame}[fragile]{İçidekiler}
    \tableofcontents
\end{frame}
%%%%%%%%%%%%%%%%%%%%%%%%%%%%%%%%%%%%%%%%%%%%%%%%%%%%%%%%%%%%%%%%%%%%%%%%%%%%%%%%
\section{Diziler yeniden}

\begin{frame}[fragile]{Dizi tanımlama}
    \begin{lstlisting}
        int len=8176*1024-10*1024;
        printf("len:%d\n",len);
        char dizi[len];
        int i=0;
        for(i=0;i<len;i++)
        {
            dizi[i]='x';
        }\end{lstlisting}
    ve
    \begin{lstlisting}
        int len=8176*1024;
        printf("len:%d\n",len);
        char dizi[len];
        int i=0;
        for(i=0;i<len;i++)
        {
            dizi[i]='x';
        }\end{lstlisting}
\end{frame}
%%%%%%%%%%%%%%%%%%%%%%%%%%%%%%%%%%%%%%%%%%%%%%%%%%%%%%%%%%%%%%%%%%%%%%%%%%%%%%%%
\begin{frame}[fragile]{Pointer(İşaretçi)}
    Bellek adresi için kullanılan veri tipine \textbf{pointer} denir.
    \begin{lstlisting}
    veri_tipi* degisken;    
    veri_tipi* degisken[];\end{lstlisting}
    ile tanımlanabilirler. Aynı zamanda
    \begin{lstlisting}
        veri_tipi degisken;    
        veri_tipi* adres=&degisken;\end{lstlisting}
    ile de kullanılabilir.
\end{frame}
%%%%%%%%%%%%%%%%%%%%%%%%%%%%%%%%%%%%%%%%%%%%%%%%%%%%%%%%%%%%%%%%%%%%%%%%%%%%%%%%
\end{document}